\documentclass[11pt,a4paper]{article}

% ---------- Packages ----------
\usepackage[margin=2.4cm]{geometry}
\usepackage{xcolor}
\usepackage{titlesec}
\usepackage{enumitem}
\usepackage{tabularx}
\usepackage{booktabs}
\usepackage{parskip}
\usepackage[colorlinks=true, linkcolor=violetaccent, urlcolor=violetaccent, citecolor=violetaccent]{hyperref}
\usepackage{fancyhdr}
\usepackage{lastpage}
\usepackage{array}

% ---------- Colors ----------
\definecolor{violetaccent}{HTML}{5A4FCF}
\definecolor{softgray}{HTML}{6B6B6B}
\definecolor{cardbg}{HTML}{F4F3FB}
\definecolor{darkbg}{HTML}{111111}

% ---------- Section styling ----------
\titleformat{\section}
  {\normalfont\Large\bfseries\color{darkbg}}
  {\thesection}{0.6em}{}
\titleformat{\subsection}
  {\normalfont\large\bfseries\color{violetaccent}}
  {\thesubsection}{0.6em}{}
\titlespacing*{\section}{0pt}{1.4em}{0.8em}
\titlespacing*{\subsection}{0pt}{1.1em}{0.5em}

% ---------- Header/Footer ----------
\pagestyle{fancy}
\fancyhf{}
\renewcommand{\headrulewidth}{0.4pt}
\fancyhead[L]{\small\color{softgray}Listening Rhythm --- Case Study}
\fancyhead[R]{\small\color{softgray}\thepage\ / \pageref{LastPage}}

% ---------- Custom "note card" environment ----------
\newenvironment{notecard}
  {\begin{quote}\color{softgray}\itshape}
  {\end{quote}}

\newcommand{\eyebrow}[1]{{\small\color{violetaccent}\textbf{\MakeUppercase{#1}}}\par\vspace{2pt}}

% ==========================================================
\begin{document}

% ---------- Title block ----------
\begin{center}
  {\eyebrow{Concept case study --- not an official Spotify product}}
  {\Huge\bfseries Listening Rhythm}\\[6pt]
  {\large\color{softgray} A UX/Product Case Study: Listener Burnout Detection for Spotify}\\[10pt]
  {\normalsize\bfseries Aarushi Singh}
\end{center}
\vspace{1.2em}
\hrule
\vspace{1.5em}

% ==========================================================
\section*{1. Executive Summary}
No --- or maybe one or two --- published case studies currently document Spotify detecting \textbf{listener burnout} specifically. The closest real research is the University of Bristol's \emph{Mood Music} study, which linked Spotify streaming history to depression, anxiety, and wellbeing scores across 171 students --- but nothing brings this out as an in-app feature.

This case study is meant to fill that gap. It proposes \textbf{Listening Rhythm}, a passive, opt-in wellbeing signal built entirely from listening \emph{behavior} (not lyrics analysis, not self-report, not biometric data) that surfaces gentle, non-diagnostic nudges when a user's patterns resemble burnout: late-night binge sessions, rising skip rates, narrowing genre diversity, and looping the same tracks --- often ones that are low in energy or fall under a slow, sad genre.

The design is built around one single hard constraint: \textbf{a music app should never play therapist --- it should be there to make a user feel light-hearted.} Every decision below --- copy, visual language, escalation path --- is shaped by that constraint.

% ==========================================================
\section*{2. Problem Framing}

\subsection*{Why this belongs on a music platform}
People already use music as an unconscious mood regulator. Research on stress recovery has shown that self-selected music listening after stressful events groups into different patterns --- calmer, sadder, more minor-key tracks --- that differ measurably from baseline listening. Spotify already has this signal for every user, but beyond recommendations, it doesn't do anything with it.

\subsection*{Why ``burnout,'' specifically}
Burnout is behavioral before it's clinical. It shows up as:

\begin{itemize}[leftmargin=1.4em]
  \item \textbf{Escape-listening} --- long, late-night sessions with minimal engagement (no skips, no searches, just letting it run)
  \item \textbf{Numbing loops} --- repeating the same 3--5 tracks for days
  \item \textbf{Energy collapse} --- a sustained drop in tempo/energy/valence of chosen tracks
  \item \textbf{Fragmented attention} --- rising skip rate, shorter session length, more genre-switching
\end{itemize}

None of these prove burnout on their own. But as a \textbf{pattern}, they're a legitimate, non-invasive signal --- closer to how a smartwatch flags irregular sleep than to a diagnosis.

\subsection*{The trap to avoid}
The obvious failure mode is a feature that says ``You seem depressed'' --- presumptuous, clinical, creepy, and panic-inducing at platform scale. The design must earn trust by staying \textbf{descriptive, not diagnostic} --- reflecting listening patterns back to the user, never naming a mental state for them.

% ==========================================================
\section*{3. Research Grounding}

\begin{table}[h!]
\renewcommand{\arraystretch}{1.4}
\begin{tabularx}{\textwidth}{@{}p{4.2cm}X@{}}
\toprule
\textbf{Source} & \textbf{Relevance} \\
\midrule
\emph{Mood Music} (Univ.\ of Bristol; EMA + Spotify streaming history + PHQ-9/GAD-7/SWEMWBS, $n=171$) &
Confirms streaming metadata correlates with real-time mood; supports using listening data as a \emph{proxy signal}, not a diagnostic tool \\
Stress-recovery listening studies &
Shows audio-feature shifts (energy, valence, tempo) are a measurable, real phenomenon post-stress --- the technical basis for a ``rhythm'' signal \\
Public playlist analysis (search terms like ``burnout,'' ``overthinking,'' ``destress'') &
Confirms users already self-label listening by emotional state, suggesting language like ``decompress'' and ``reset'' will feel native rather than clinical \\
\bottomrule
\end{tabularx}
\end{table}

\noindent\textbf{Design implication:} the feature should read signals people are already using implicitly (they already build ``de-stress'' playlists) rather than inventing a new clinical term.

% ==========================================================
\section*{4. Persona}
\textbf{Priya, 29, product designer.}
A Premium user who listens to music for 3+ hours a day. For the last few weeks she's been going to bed later, her ``Discover Weekly'' engagement has dropped, and she's been replaying the same lo-fi playlist every night instead of skipping through it like she used to. She hasn't consciously noticed any of this. She would be irritated by a pop-up asking ``Are you okay?'' --- but might appreciate a quiet, optional insight, on her own terms, that she can dismiss in one tap.

Priya is the target: \textbf{not someone in crisis}, but someone in a slow behavioral drift who has no reason to notice it herself --- or who is somehow, knowingly, avoiding noticing this shift in her own pattern.

% ==========================================================
\section*{5. Ethical Guardrails}
\emph{(non-negotiable constraints on the design)}

\begin{enumerate}[leftmargin=1.6em]
  \item \textbf{Opt-in only}, off by default. No passive shipping of a ``we noticed you seem stressed'' feature to 600M+ users without consent --- this could create panic among users.
  \item \textbf{Never name a clinical state.} No ``depression,'' ``burnout,'' or ``anxiety'' labels attached to the user. Only describe \emph{behavior} (``your late-night sessions are up'') and let the user draw their own conclusion --- this may also help them realize what they've been doing over the past few days.
  \item \textbf{No push notifications for this.} It lives in-app, checked only when the user opens the app --- never an interruption.
  \item \textbf{One-tap exit at every step.} Dismiss, mute for 30 days, or turn off entirely --- always visible, never hidden in settings or behind multiple three-dot menus.
  \item \textbf{A support resource is offered, never forced.} If a user chooses to dig deeper, options like grounding exercises or, if they ask, professional resources are one tap away --- but never inserted automatically or repeatedly. Let users be their own guardians.
  \item \textbf{No data leaves the feature.} Not shared with advertisers, not used to target ads --- stated explicitly in the opt-in screen, since trust is the entire foundation for this feature.
\end{enumerate}

% ==========================================================
\section*{6. The Feature: ``Listening Rhythm''}

A single visual metaphor carries the whole feature: \textbf{your listening as a rhythm/pulse.} Steady, evenly-spaced pulses = a stable pattern. Clustered, late, erratic pulses = a shifting one. It's the same metaphor on every screen, so the feature never needs to explain itself twice.

\subsection*{Signals tracked}
\emph{(all passive, all already-collected metadata)}
\begin{itemize}[leftmargin=1.4em]
  \item Session start times (late-night clustering)
  \item Skip rate trend (7-day rolling average vs.\ personal baseline)
  \item Repeat-track ratio (same track/playlist looped vs.\ personal baseline)
  \item Audio-feature drift (average energy/valence/tempo of chosen tracks vs.\ personal baseline)
\end{itemize}

Everything is \textbf{relative to the user's own baseline} --- never compared to other users. This avoids both the creepiness of population-level profiling and the inaccuracy of one-size-fits-all thresholds.

\subsection*{User flow}
\begin{notecard}
Discover tab (existing)
$\rightarrow$ Opt-in card, shown once
\quad ``Not now'' $\rightarrow$ never shown again unless user seeks it out
\quad ``Try it'' $\rightarrow$ passive tracking begins (invisible, no UI change)
$\rightarrow$ Pattern shift detected over rolling window
$\rightarrow$ Quiet insight card appears at bottom of Home feed (not a notification)
\quad ``Not now'' $\rightarrow$ dismissed, resurfaces no sooner than +14 days
\quad ``Tell me more'' $\rightarrow$ Weekly Rhythm report
\quad (ignored) $\rightarrow$ disappears after this session, no nagging
$\rightarrow$ Weekly Rhythm report (opt-in deep dive)
\quad ``Reset playlist'' (mood-lift suggestions)
\quad ``Just data, thanks'' (closes, no action)
\quad ``I'd like support resources'' $\rightarrow$ optional, user-initiated only
\end{notecard}

% ==========================================================
\section*{7. Key Screens}
See the accompanying interactive mockup for the four core screens:
\begin{enumerate}[leftmargin=1.6em]
  \item \textbf{Opt-in card} --- plain-language consent, states exactly what is and isn't tracked
  \item \textbf{Home feed insight card} --- one line, one visual, two buttons, easily ignorable
  \item \textbf{Weekly Rhythm report} --- the full behavioral picture, framed as data, not diagnosis
  \item \textbf{Reset playlist} --- the one actionable output, turning every insight into something useful immediately
\end{enumerate}

% ==========================================================
\section*{8. Success Metrics}

\begin{table}[h!]
\renewcommand{\arraystretch}{1.4}
\begin{tabularx}{\textwidth}{@{}p{5.2cm}X@{}}
\toprule
\textbf{Metric} & \textbf{What it tells us} \\
\midrule
Opt-in rate & Whether the concept is trusted enough to try \\
Dismiss-without-return rate & Whether the insight card feels like nagging \\
``Reset playlist'' save rate & Whether the feature offers real value, not just surveillance \\
Opt-out rate after first insight & The clearest signal of whether tone/trust failed \\
Session length \emph{after} engaging with a Reset playlist & Whether the feature actually helps, however loosely defined \\
\bottomrule
\end{tabularx}
\end{table}

We deliberately exclude clinical detection ``accuracy'' as a success metric. The feature is intended to help users reflect on their listening patterns and overall wellbeing, not to provide medical assessments or burnout diagnoses.

% ==========================================================
\section*{9. Risks \& Open Questions}
\begin{itemize}[leftmargin=1.4em]
  \item \textbf{False positives at scale}: a user who has simply started a new night-shift job will trigger the same signals as someone withdrawing. The copy must stay behavioral (``your listening times have shifted'') rather than interpretive (``you seem burnt out'') so a wrong guess doesn't feel invasive.
  \item \textbf{Regulatory exposure}: even non-clinical wellbeing inference from behavioral data increasingly draws scrutiny (EU AI Act ``emotion inference'' provisions, for one). This would need real legal review before ship.
  \item \textbf{Trust is the whole product.} If this ever feels like it's profiling users for ad targeting --- even by rumor --- the feature dies overnight. The ``no data leaves this feature'' guarantee has to be real and independently auditable, not just a line of copy.
\end{itemize}

% ==========================================================
\section*{10. Why This Belongs in a Portfolio}
This case study deliberately treats the hardest part of the problem --- ethics --- as the core design challenge, rather than an afterthought. A simpler approach might be to recommend a mood-boosting playlist whenever a user feels sad or emotional. But the real challenge here is to design a feature that genuinely supports users instead of making Spotify feel like it's making assumptions about someone's mental health. Handling that balance is what makes this a meaningful product.

\end{document}
